\section{Votum}
\label{sec:votum}

\subsection{Die Entscheidung}

\vd{\textbf{Nicht kaufen zu \je{\Kurs}{EUR}. Beobachten, und bei einem Kurs unter
\je{\SchwelleMEDZehnLang}{EUR} erneut rechnen.} Wer die Aktie h\"alt, hat keinen Grund zu
verkaufen: Das Gesch\"aft ist stetig, schuldenfrei und w\"achst; das Votum richtet sich
gegen den Preis, nicht gegen das Unternehmen.}

\vd{Der Horizont dieses Votums ist zehn Jahre, weil auf f\"unf Jahre jede Anlage dieser Art
am Endwert h\"angt und auf f\"unfzehn Jahre die konstante Fortschreibung des Zuflusses
nichts mehr aussagt.}

\subsection{Der Preis, zu dem die Antwort anders lautete}

\schwellentabelle

\vd{Bei der H\"urde von \Pz{\Huerde} liegt der Schwellenkurs auf zehn Jahre im Basisfall bei
\je{\SchwelleMEDZehn}{EUR}, bei der vorsichtigeren H\"urde von \Pz{\HuerdeLang} bei
\je{\SchwelleMEDZehnLang}{EUR}. Der heutige Kurs ist das \VielfachesMED-fache des ersten
Werts. Selbst im optimistischen Szenario und bei der milderen H\"urde erreicht der
Schwellenkurs nur \je{\SchwelleOPTZehnLang}{EUR}.}

\subsection{Was der heutige Kurs voraussetzt}

\umkehrtabellen

Zum heutigen Kurs erreicht die Anlage die H\"urde auf zehn Jahre nur, wenn der freie
Zufluss um \WachstumMEDZehn{} p.\,a.\ w\"achst oder wenn das Unternehmen am Ende des
Horizonts das \EndwertMEDZehn-fache seines heutigen Buchwerts wert ist.

Diesen beiden Zahlen stehen drei belegte Gr\"o{\ss}en gegen\"uber. Erstens die eigene
Reihe: Der freie Zufluss ist zwischen \FcfJahre{} nicht gewachsen, sein h\"ochster Wert
stammt aus dem Jahr \FcfMaxJahr. Zweitens die Ziele des Unternehmens: Mission 30
impliziert ein EBITDA-Wachstum von \MissionImplizit{} p.\,a.\ und damit weniger als ein
Viertel des verlangten Werts. Drittens die eigene Bewertungsgeschichte: Das KGV von
\KgvHeute{} liegt \"uber dem Median von \KgvMedian{} und damit \"uber
\Pz{\KgvRang} der Jahre seit 2020, bleibt aber innerhalb der eigenen Spanne.

\vd{Die Zahl, die das Votum tr\"agt, ist \WachstumMEDZehn{} gegen \MissionImplizit. Der
Kurs verlangt ein Vielfaches dessen, was das Unternehmen sich selbst vornimmt, und ein
Vielfaches dessen, was es in sieben Jahren erreicht hat. Das ist keine Aussage \"uber die
Qualit\"at von GEA, sondern \"uber den Abstand zwischen Preis und belegtem Zahlungsstrom.}

\subsection{Was das Votum \"andern w\"urde}

\vd{Erstens: ein Gesch\"aftsjahr, in dem der freie Zufluss deutlich \"uber
\Mio{\FcfMax}{EUR} steigt, also der Punkt, an dem sich die erh\"ohten Investitionen der
Jahre 2022 bis 2025 im Zufluss zeigen. Der Auftragsbestand von
\Mio{\HjAuftragsbestand}{EUR} und das Halbjahreswachstum von \HjAuftragWachstum{} machen
das f\"ur 2026 oder 2027 m\"oglich; dieser Bericht rechnet nicht damit, sondern pr\"uft es
nach.}

\vd{Zweitens: ein Kurs nahe \je{\SchwelleMEDZehnLang}{EUR}. Drittens: ein belegter
Marktanteil, der die im Bericht behauptete F\"uhrungsposition st\"utzt und damit einen
Endwert \"uber dem Buchwert rechtfertigt.}

\subsection{Was gegen dieses Votum spricht}

\vd{Der st\"arkste Einwand ist der Endwert. Die Rechnung gibt GEA nach zehn Jahren nur
seinen heutigen Buchwert von \je{\BuchwertJeAktie}{EUR} je Aktie. F\"ur einen Maschinenbauer
mit \Pz{\Serviceanteil} Serviceanteil, \Pz{\Roce} ROCE und Nettoliquidit\"at ist das
streng; ein Unternehmen, das sein eingesetztes Kapital derart verzinst, wird am Markt nie
zum Buchwert gehandelt. Wer stattdessen den heutigen Multiplikator als Endwert ansetzt,
kommt zu einem h\"oheren Schwellenkurs.}

\vd{Der zweite Einwand sind die Organk\"aufe. \DdKaeufe{} K\"aufe von \DdPersonen{}
Vorstandsmitgliedern ohne einen einzigen Verkauf sind das deutlichste Insidersignal, das
dieser Bericht gefunden hat, und die K\"aufer kennen die Auftragslage besser als jede
Reihe. Ihre K\"aufe lagen allerdings zwischen \je{\DdPreisVon}{EUR} und
\je{\DdPreisBis}{EUR}, also unter dem heutigen Kurs; sie sprechen f\"ur das Unternehmen
und f\"ur jenen Preis, nicht f\"ur diesen.}

\vd{Der dritte Einwand ist die Prognosetreue: \PrognoseTreffer{} von
\PrognoseQuantifiziert{} Prognosen erreicht oder \"ubertroffen. Wer daraus schlie{\ss}t,
dass Mission 30 eher \"ubertroffen als verfehlt wird, muss die \MissionImplizit{} als
Untergrenze lesen. Auch dann bleibt der Abstand zu \WachstumMEDZehn{} gro{\ss}.}
